% =====================================================================
% Real-Time FPGA-Based Pulse Simulator for Calorimeter Readout
% Esqueleto de artigo de periódico — classe IEEEtran.
% REGRAS (revisão do Luciano, 2026-07-17):
%   (1) citações: QUALQUER tipo de fonte (periódico, conferência,
%       preprint, datasheet/módulo comercial — ex.: CAEN), desde que
%       o mais RECENTE possível e em INGLÊS;
%   (2) NENHUMA menção a ATLAS/TileCal/LHC — enquadrar como
%       "nuclear instrumentation" genérica.
% Compilar: pdflatex -> bibtex -> pdflatex x2
% =====================================================================
\documentclass[journal]{IEEEtran}

\usepackage{graphicx}
\usepackage{amsmath}
\usepackage{cite}
\usepackage{xcolor}
\usepackage{url}

% Notas TODO visíveis no PDF — remover antes da submissão
\newcommand{\todo}[1]{\textcolor{red!70!black}{\textit{#1}}}

\begin{document}

\title{Real-Time FPGA-Based Pulse Simulator for Calorimeter Readout Electronics in Nuclear Instrumentation}

\author{Fábio~C.~Luna,
        Thiago~C.~A.~Paschoalin,
        Tiago~M.~Quirino,
        and~Luciano~M.~de~Andrade~Filho
% TODO: confirmar lista e ordem de autores, afiliações e ORCIDs
\thanks{F. C. Luna and L. M. de Andrade Filho are with the Electrical Engineering Department, Federal University of Juiz de Fora, Juiz de Fora, Brazil.}%
\thanks{T. C. A. Paschoalin is with the Electroelectronic Department, CEFET-MG, Leopoldina, Brazil.}%
\thanks{T. M. Quirino is with the Electrical Engineering Department, Rio de Janeiro State University, Rio de Janeiro, Brazil.}}

\maketitle

\begin{abstract}
\todo{($\sim$200 palavras) Contexto: processamento de pulsos em instrumentação nuclear / calorimetria de altas energias. Problema: validar técnicas de reconstrução sem acesso ao experimento (custo, acesso restrito, radiação). Proposta: simulador de pulsos em tempo real em FPGA a 40 MHz, com três blocos — RNG, atribuição de energia das colisões, shaper. Resultados principais: recursos $< X\%$, $f_{max}$, fidelidade do espectro e do pulso. Atenção: sem citar o experimento/detector específico.}
\end{abstract}

\begin{IEEEkeywords}
FPGA, real-time simulation, pulse simulator, nuclear instrumentation, calorimetry, random number generation, digital shaper.
\end{IEEEkeywords}

% =====================================================================
\section{Introduction}

\todo{Instrumentação nuclear e calorímetros para medição de energia de partículas em experimentos de altas energias (genérico, sem citar o experimento).}

\todo{Problema: desenvolver e validar técnicas de processamento diretamente no setup experimental é impraticável (custo, acesso, radiação); simuladores de pulso em tempo real como ambiente controlado.}

\todo{Vantagem do FPGA vs.\ simulação offline: latência, paralelismo, hardware-in-the-loop.}

\todo{Breve menção aos trabalhos do grupo \cite{quirino2025,paschoalin2026} e apontar para a seção de estado da arte (Section II).}

\todo{Contribuição: simulador completo em três estágios — RNG, atribuição de energia das colisões, shaper — gerando pulsos a 40 MHz, compatível com a taxa de eventos de experimentos de colisão modernos. Organização do texto.}

% =====================================================================
\section{Real-Time Pulse Generation: State of the Art}
% Nova seção pedida na revisão: responder (i) o que existe de
% simuladores/geradores de pulso em tempo real para física nuclear,
% (ii) com que arquitetura cada um trabalha, (iii) em que nosso método
% difere.

\todo{Síntese digital de formas de pulso como origem da linha \cite{jordanov1994,regadio2014}.}

\todo{Geradores em FPGA evento a evento: multicanal Poisson \cite{arkani2013}, amplitude com distribuição prescrita \cite{ponikvar2018}, formas exponenciais/gaussianas configuráveis \cite{garciaduran2019}, emulador Poisson por ensaios de Bernoulli \cite{ieee_neutron_emulator}.}

\todo{Rota ``detector virtual''/software: emulador com pile-up \cite{pechousek2016}, Monte Carlo \cite{lisangang2021,ma2024}, nicho de nêutrons/reatores \cite{wen2022}.}

\todo{Equipamento comercial: família CAEN Digital Detector Emulator (DT5800/DT5810/DT5810B/DT4800) \cite{caen_dt5810b}, tecnologia publicada em \cite{abba2014}. Posicionamento: ``commercially available but costly'', arquitetura proprietária, orientado a espectroscopia evento a evento — NÃO inventar preço.}

\todo{Competidor mais recente e próximo: emulação de detector SiPM em FPGA em tempo real \cite{carsi2026} (preprint jul/2026, submetido ao IEEE TNS — atualizar citação se sair a versão de periódico).}

\todo{FIO CONDUTOR (argumento central): os geradores existentes são evento a evento, de taxa baixa, com formas simples (exponencial/gaussiana) e estatística Poisson; nenhum modela a cadeia completa de leitura (shaper IIR de cauda longa) em regime de colisão a 40 MHz, com ocupância e espectro configuráveis por célula. No nosso simulador, pile-up e desvio de linha de base não são ``modos'' do gerador: emergem naturalmente da superposição dos pulsos na taxa de colisão.}

% =====================================================================
\section{Calorimeter Readout Chain Model}

\todo{Cadeia de leitura com cintilador + PMT como \emph{exemplo} (não como \emph{a} cadeia): cintilador $\rightarrow$ PMT $\rightarrow$ condicionamento (CSP + shaper) $\rightarrow$ ADC a 40 MHz \cite{quirino2025}. Citar ao menos uma alternativa de geração do sinal de prompt (ex.\ ionização de gás: câmaras de ionização, contadores proporcionais) e deixar explícito que o simulador requer apenas o pulso característico $h[n]$ do canal — no modelo convolucional, só $h[n]$ muda, a arquitetura é a mesma.}

\todo{Modelo convolucional: deposições de energia como impulsos discretos $s[n]$, dinâmica do canal concentrada no pulso característico $h[n]$:}
\begin{equation}
  r[n] = h[n] \ast s[n]
\end{equation}

\todo{Premissa que justifica o desenho do simulador (não deixar como frase solta — fazer a ponte causal): ``since the front-end electronics are assumed identical across the channels of a given system, a single characteristic pulse $h[n]$ describes the system; cell-to-cell differences reduce to occupancy and energy spectrum, precisely the parameters the simulator exposes as runtime-configurable''. Alternativa: mover para o início da Section IV (arquitetura) como motivação da parametrização. Figura: diagrama de blocos da cadeia de leitura.}

% =====================================================================
\section{Simulator Architecture and Implementation}

\todo{Visão geral do pipeline em três blocos + figura do diagrama geral.}

\subsection{Random Number Generation}

\todo{Requisito: sequências independentes e não correlacionadas para a decisão de evento e o sorteio de energia \cite{coates1988}.}

\todo{Algoritmo adotado (ex.\ abordagem multi-memória \cite{paschoalin2026}); comparação com LFSR, Mersenne Twister e xorshift \cite{matsumoto1998,marsaglia2003,thomas2007,li2014}.}

\todo{Implementação em FPGA: recursos, período, testes estatísticos.}

\subsection{Collision Energy Assignment}

\todo{Ocupação: comparação do valor aleatório uniforme com a taxa de ocupação da célula $\rightarrow$ sequência binária de eventos.}

\todo{Amplitude: CDF inversa do espectro de energia da célula, discretizada em LUT (128--4096 entradas) em memória.}

\todo{Parâmetros configuráveis em tempo real (ocupação, offset, espectro). Figura: espectro de energia alvo vs.\ gerado.}

\subsection{Shaper}

\todo{Estágio de conformação: CSP + PSC (CR-4RC), resposta com cauda longa \cite{quirino2025}.}

\todo{FIR inviável ($\sim 10^4$ coeficientes) $\rightarrow$ shaper IIR: discretização por invariância ao impulso, decomposição em seções paralelas de 1\textsuperscript{a}/2\textsuperscript{a} ordem, aritmética de ponto fixo \cite{jordanov1994,regadio2014}.}

\todo{Figura: resposta ao impulso ideal vs.\ implementada; efeito de pile-up e desvio de linha de base.}

% =====================================================================
\section{Results}

\todo{Recursos do FPGA (ALMs/LUTs, DSPs, memória), frequência máxima.}

\todo{Validação: comparação com referência ideal (erro de quantização), fidelidade do espectro e da forma de pulso.}

\todo{Operação a 40 MHz com parâmetros ajustados em tempo real. Tabelas e figuras de resultados.}

% =====================================================================
\section{Conclusions}

\todo{Síntese das contribuições, limitações e trabalhos futuros (integração com algoritmos de reconstrução, múltiplos canais, cenários de alta taxa de eventos).}

\section*{Acknowledgments}

\todo{Agências de fomento (CNPq, FAPEMIG, FAPERJ, CAPES\ldots).}

\bibliographystyle{IEEEtran}
\bibliography{references}

\end{document}
